\section{Obiekty uniwersalne: początkowe, końcowe, zerowe}\label{sec:uniwersalne}

\begin{definicja}[Obiekt początkowy i końcowy]\label{def:poczatkowy}
Obiekt $I$ kategorii $\mathcal C$ jest \emph{początkowy}, jeśli dla każdego $A\in\Ob(\mathcal C)$ istnieje dokładnie jeden morfizm $I\to A$. Obiekt $T$ jest \emph{końcowy}, jeśli dla każdego $A\in\Ob(\mathcal C)$ istnieje dokładnie jeden morfizm $A\to T$. Obiekt będący jednocześnie początkowym i końcowym nazywamy \emph{obiektem zerowym}.
\end{definicja}

\begin{przyklad}[Zbiory, grupy, przestrzenie]
W $\Set$ zbiór pusty $\varnothing$ jest początkowy, ponieważ istnieje dokładnie jedna funkcja $\varnothing\to A$ dla każdego zbioru $A$. Dowolny singleton $\{*\}$ jest końcowy: jedyna możliwa funkcja przypisuje każdemu elementowi $A$ punkt $*$ (również dla $A=\varnothing$). W $\Grp$ grupa trywialna $\{e\}$ jest jednocześnie początkowa i końcowa, czyli zerowa: homomorfizm z niej musi wysłać $e$ na element neutralny, a do niej istnieje tylko funkcja stała, która jest homomorfizmem. W $\Vect_k$ analogicznie obiektem zerowym jest przestrzeń $\{0\}$, ponieważ odwzorowania liniowe z niej i do niej są jednoznacznie wyznaczone. W $\Top$ obiektem początkowym jest przestrzeń pusta, a końcowym dowolna przestrzeń jednoelementowa: jedyne funkcje o odpowiednich typach są ciągłe, gdyż przeciwobrazy zbiorów otwartych są puste lub równe całej dziedzinie.
\end{przyklad}

\begin{przyklad}[Minimum i maksimum]
W kategorii częściowego porządku $(P,\leq)$ obiekt początkowy to najmniejszy element $P$, jeśli istnieje; obiekt końcowy to największy element. Na przykład w $(\mathcal P(S),\subseteq)$ są to odpowiednio $\varnothing$ i $S$.
\end{przyklad}

\begin{twierdzenie}[Jednoznaczność z dokładnością do jedynego izomorfizmu]\label{tw:unique-universal}
Jeśli $I,J$ są obiektami początkowymi tej samej kategorii, to istnieje dokładnie jeden izomorfizm $I\to J$. Analogicznie dla dwóch obiektów końcowych.
\end{twierdzenie}
\begin{proof}
Z początkowości $I$ istnieje jedyny $f:I\to J$, a z początkowości $J$ jedyny $g:J\to I$. Złożenie $g\circ f:I\to I$ musi być równe $\id_I$, ponieważ istnieje dokładnie jeden morfizm $I\to I$; podobnie $f\circ g=\id_J$. Jedyność izomorfizmu wynika z jedyności dowolnego morfizmu $I\to J$. Jeżeli $T,U$ są końcowe, końcowość $U$ daje jedyny morfizm $a:T\to U$, a końcowość $T$ jedyny $b:U\to T$. Ponieważ $b\circ a$ i $\id_T$ są morfizmami do $T$, są równe; analogicznie $a\circ b=\id_U$. Jedyność izomorfizmu wynika z końcowości $U$. W szczególności dwa obiekty zerowe są połączone jedynym izomorfizmem.
\end{proof}

\begin{uwaga}
Wyrażenie \emph{jednoznaczny z dokładnością do jedynego izomorfizmu} jest mocniejsze od zwykłego \emph{jednoznaczny z dokładnością do izomorfizmu}: nie tylko obiekty muszą być izomorficzne, lecz także odpowiedni izomorfizm jest wymuszony przez daną własność uniwersalną. Ta idea wróci przy iloczynach, granicach i sprzężeniach.
\end{uwaga}

\subsection{Dalsze przykłady i ograniczenia}

\begin{przyklad}[Pierścienie i moduły]\label{cat:ex:initial-ring}
W kategorii pierścieni przemiennych z jedynką i homomorfizmów zachowujących jedynkę obiektem początkowym jest $\mathbb Z$. Jedyny homomorfizm do pierścienia $R$ ma postać $n\mapsto n1_R$: addytywność i zachowanie jedynki wymuszają ten wzór, który zachowuje także mnożenie, gdyż $(n1_R)(m1_R)=(nm)1_R$. Przyjmujemy konwencję z rozdziału 9 skryptu \emph{Algebra abstrakcyjna}, dopuszczającą pierścień zerowy, czyli $0=1$. Jest on obiektem końcowym: jedyna funkcja $R\to\{0\}$ zachowuje wszystkie działania i jedynkę, która w przeciwdziedzinie jest zerem. Nie istnieje obiekt zerowy, gdyż $\mathbb Z$ nie jest izomorficzny z pierścieniem zerowym.

Gdyby wykluczyć pierścień zerowy, nadal istniałby obiekt początkowy $\mathbb Z$, lecz nie istniałby obiekt końcowy. Istotnie, homomorfizmy z $\mathbb F_2$ i $\mathbb F_3$ do hipotetycznego obiektu końcowego $T$ wymuszałyby jednocześnie $2\cdot1_T=0$ i $3\cdot1_T=0$, a po odjęciu $1_T=0$, sprzecznie z przyjętą konwencją.

Dla ustalonego pierścienia $R$ z jedynką kategoria lewych unitarnych $R$-modułów ma obiekt zerowy $\{0\}$. Każdy homomorfizm z modułu zerowego lub do niego jest zerowy, a taki homomorfizm zawsze istnieje. Moduł ten jest również modułem wolnym na pustym zbiorze generatorów (por. rozdział 16 skryptu \emph{Algebra abstrakcyjna}).
\end{przyklad}

\begin{przyklad}[Przestrzenie z wyróżnionym punktem i kategoria relacji]\label{cat:ex:pointed-zero}
W kategorii $\Top_*$ przestrzeni z wyróżnionym punktem singleton z jego jedynym punktem wyróżnionym jest zerowy. Jedyna dopuszczalna strzałka ze singletonu do $(X,x_0)$ wybiera $x_0$ i jest ciągła; jedyna strzałka w przeciwną stronę jest stała i zachowuje wyróżniony punkt. Zbiór pusty nie jest obiektem tej kategorii.

W kategorii $\Rel$ relacji z rozdziału~\ref{sec:kategorie} zbiór pusty jest zerowy, ponieważ zarówno $\varnothing\times A$, jak i $A\times\varnothing$ są puste i każdy ma dokładnie jeden podzbiór. Widzimy, że ten sam zbiór może być tylko początkowy w $\Set$, ale zerowy w $\Rel$, zależnie od tego, jakie morfizmy dopuszczamy.
\end{przyklad}

\begin{przyklad}[Kategorie bez obiektu początkowego lub końcowego]
W kategorii odpowiadającej grupie nietrywialnej z jednym obiektem jedyny obiekt nie jest ani początkowy, ani końcowy: ma więcej niż jeden endomorfizm. W kategorii zbiorów i funkcji różnowartościowych zbiór pusty jest początkowy, ale obiekt końcowy nie istnieje. Gdyby $T$ był końcowy, iniekcje ze singletonu do $T$ odpowiadałyby jego elementom; zatem $T$ miałby dokładnie jeden element. Nie istnieje jednak iniekcja ze zbioru dwuelementowego do singletonu.
\end{przyklad}

\begin{stwierdzenie}[Własności obiektów początkowych i końcowych]\label{cat:prop:initial-terminal-arrows}
Każdy morfizm \emph{wychodzący} z obiektu końcowego jest monomorfizmem rozszczepionym, a każdy morfizm \emph{wchodzący} do obiektu początkowego jest epimorfizmem rozszczepionym.
\end{stwierdzenie}
\begin{proof}
Niech $T$ będzie końcowy i $f:T\to A$. Jedyny morfizm $r:A\to T$ spełnia $r\circ f=\id_T$, ponieważ $T$ ma dokładnie jeden endomorfizm. Jeśli $I$ jest początkowy i $g:A\to I$, jedyny morfizm $s:I\to A$ spełnia $g\circ s=\id_I$. Otrzymujemy odpowiednie rozszczepienia z rozdziału~\ref{sec:morfizmy}. Dla kontrastu morfizm $\varnothing\to\{*\}$ w $\Set$ nie jest epimorfizmem, mimo że jego przeciwdziedzina jest obiektem końcowym.
\end{proof}

\begin{definicja}[Ścisły obiekt początkowy]\label{cat:def:strict-initial}
Obiekt początkowy $I$ jest \emph{ścisły}, jeżeli każdy morfizm $A\to I$ jest izomorfizmem.
\end{definicja}

\begin{przyklad}
W $\Set$ i $\Top$ obiekt pusty jest ścisły: istnienie funkcji $A\to\varnothing$ wymusza pustą dziedzinę, a jedyne odwzorowanie między pustymi obiektami jest izomorfizmem. W $\Grp$ obiekt początkowy nie jest ścisły, ponieważ homomorfizm z dowolnej nietrywialnej grupy do grupy trywialnej nie jest izomorfizmem.
\end{przyklad}

\subsection{Morfizmy zerowe}

\begin{definicja}[Morfizm zerowy]\label{cat:def:zero-morphism}
Jeżeli kategoria ma obiekt zerowy $O$, to dla dowolnych obiektów $A,B$ definiujemy morfizm zerowy $0_{A,B}:A\to B$ jako złożenie
\[
\begin{CD}
A @>>> O @>>> B.
\end{CD}
\]
Pierwsza strzałka jest jedyna dzięki końcowości $O$, a druga dzięki jego początkowości.
\end{definicja}

\begin{stwierdzenie}[Niezależność wyboru i pochłanianie przez zero]\label{cat:prop:zero-absorption}
Morfizm zerowy nie zależy od wyboru obiektu zerowego. Dla $u:X\to A$ oraz $v:B\to Y$ zachodzą
\[
0_{A,B}\circ u=0_{X,B},\qquad
v\circ0_{A,B}=0_{A,Y}.
\]
Ponadto obiekt $A$ jest zerowy wtedy i tylko wtedy, gdy $\id_A=0_{A,A}$.
\end{stwierdzenie}
\begin{proof}
Oznaczmy jedyne strzałki $A\to O$ i $O\to B$ przez $p_A$ i $j_B$. Jeżeli $O'$ jest drugim obiektem zerowym, jedyny izomorfizm $h:O\to O'$ spełnia $p'_A=h\circ p_A$ i $j'_B\circ h=j_B$, z jedyności strzałek do $O'$ i z $O$. Zatem $j'_B\circ p'_A=j_B\circ p_A$.

Ponieważ $p_A\circ u$ jest strzałką $X\to O$, mamy $p_A\circ u=p_X$. Podobnie $v\circ j_B=j_Y$. Stąd
\[
(j_B\circ p_A)\circ u=j_B\circ p_X,
\qquad v\circ(j_B\circ p_A)=j_Y\circ p_A.
\]
Jeżeli $A$ jest zerowy, ma jedyny endomorfizm, więc $\id_A=0_{A,A}$. Odwrotnie, przy tej równości dla każdego $f:A\to B$ mamy $f=f\circ\id_A=f\circ0_{A,A}=0_{A,B}$. Dla każdego $g:B\to A$ analogicznie $g=\id_A\circ g=0_{B,A}$. Istnienie morfizmów zerowych zapewnia, że oba zbiory morfizmów są jednoelementowe, więc $A$ jest zerowy.
\end{proof}

\begin{przyklad}[Postać morfizmów zerowych]
W kategorii grup $0_{G,H}$ jest homomorfizmem wysyłającym każdy element $G$ na element neutralny $H$; dla modułów jest zwykłym odwzorowaniem zerowym. W $\Top_*$ jest odwzorowaniem stałym o wartości w punkcie wyróżnionym przeciwdziedziny, a w $\Rel$ relacją pustą. Każda z tych postaci wynika ze złożenia przez odpowiedni obiekt zerowy.
\end{przyklad}

\begin{stwierdzenie}[Funktory zachowujące obiekt zerowy]\label{cat:prop:zero-functor}
Jeżeli $F:\mathcal C\to\mathcal D$ jest funktorem między kategoriami mającymi obiekty zerowe i $F$ przeprowadza obiekt zerowy $\mathcal C$ na obiekt zerowy $\mathcal D$, to $F(0_{A,B})=0_{F(A),F(B)}$.
\end{stwierdzenie}
\begin{proof}
Z funktorialności $F(j_B\circ p_A)=F(j_B)\circ F(p_A)$. Prawa strona faktoryzuje się przez obiekt zerowy $F(O)$, a obie strzałki do i z niego są jedyne. Niezależność wyboru obiektu zerowego daje tezę.
\end{proof}

\subsection{Jednoznaczność jako schemat własności uniwersalnej}

\begin{uwaga}\label{cat:rem:universal-coherence}
Początkowość i końcowość są własnościami obiektów, ale w dalszych konstrukcjach uniwersalnych trzeba pamiętać także o dodatkowych danych (np. rzutach iloczynu z rozdziału~\ref{sec:produkty}). Jedyny izomorfizm jest wtedy jedyny wśród morfizmów zgodnych z tymi danymi, niekoniecznie wśród wszystkich morfizmów między obiektami bazowymi.

Jeżeli $I,J,K$ są początkowe, a $u_{I,J}:I\to J$ oznacza jedyny morfizm, to $u_{J,K}\circ u_{I,J}=u_{I,K}$, gdyż obie strony są morfizmami z $I$ do $K$. Podobna równość zachodzi dla obiektów końcowych. Własność uniwersalna wymusza zatem nie tylko istnienie porównań, lecz także ich zgodność przy składaniu.
\end{uwaga}

